\chapter*{Integrated Synthesis and Open Problems}
\addcontentsline{toc}{chapter}{Integrated Synthesis and Open Problems}
\label{back:final-synthesis}

Version 0.11 reconciles the publication text with the current repository mathematics. The canonical numbered theorem sequence remains through SOH-G023. The integrated G024 line contains proved first- and second-order signed derivative results, the explicit computer-assisted Q dependency, and the reviewed SOH-G024-T theorem closing the full complete-monotonicity route. The reciprocal--conjugation projective chapter gives a minimal exact RH-equivalent orbit formulation.

\section*{Closed results and closed routes}
The monograph establishes the half-axis fixed geometry, exact projective coordinate identities, the centred quotient $\xi(1/2+z)=F(z^2)$, the positive Riemann-kernel representation, PF$_2$ at the declared coefficient level, the negative-inversion spectral no-go, and the G021--G023 coefficient reductions.

For
\[
H_y(q)=D_y(\sqrt q),\qquad0<|y|<\frac12,
\]
G024 proves
\[
H_y'(q)<0,\qquad H_y''(q)>0\qquad(q\ge0).
\]
The second result uses the explicit outward-interval certificate for $L''''<20L''$ on its compact trust region plus analytic continuation of the estimate to the tail.

SOH-G024-T adds a qualitatively different result. The Riemann-kernel tail and strong convexity imply
\[
\forall T>0:\quad e^{Tq}H_y(q)\to0.
\]
No nonzero completely monotone function with finite $h(0)$ can have this decay, because its Bernstein representation supplies a positive exponential lower envelope. Therefore
\[
\boxed{H_y\text{ is not completely monotone}.}
\]
The all-order Gaussian-mixture route is therefore CLOSED, not merely unfinished. The first two derivative theorems remain valid.

\section*{Exact reciprocal--conjugation minimum cut}
With
\[
u=\Omega(s)=\frac{s}{1-s},\qquad X(u)=\xi\!\left(\frac{u}{1+u}\right),
\]
the functional equation and conjugation act by
\[
\mathcal I(u)=1/u,\qquad\mathcal C(u)=\bar u.
\]
Their pointwise coincidence is
\[
\boxed{\mathcal I(u)=\mathcal C(u)\iff |u|=1\iff\Re s=\frac12.}
\]
Thus
\[
\boxed{\mathrm{RH}\iff \forall u\in Z_X:\;u^{-1}=\bar u.}
\]
The defect
\[
\Delta_{\mathrm{RC}}(u)=\left|u^{-1}-\bar u\right|^2
\]
is non-negative and vanishes exactly on the critical line. This removes the ambiguity between a self-dual point and an off-axis reciprocal two-cycle that is invisible to orbit barycentres. It does not independently prove that the defect vanishes at every zero.

\section*{Occam relational zero triad}

Chapter~\ref{chap:relational-zero-triad} adds an orthogonal theorem line with only
two declared structural axioms: normalized complementarity and a positive
relational action-defect. It proves the same local zero by three routes:
\[
J(\sigma)=\sigma
\iff
\sigma=\frac12,
\]
\[
\ln2-H_2(\sigma)=0
\iff
\sigma=\frac12,
\]
and, for a finite or countable relation family,
\[
\mathfrak S_{\rm rel}=0
\iff
\forall e:\quad F_e=0,\quad \sigma_e=\frac12.
\]
This closes the internal relational-zero theorem SOH-RZ001--SOH-RZ003. The
project defect \(\mathfrak S_{\rm rel}=0\) is not identified with the ordinary
mechanical statement \(\delta S=0\), and the project term ``Lagrange node''
denotes a zero of the declared relational vector field rather than a claim
that all astronomical Lagrange points are geometric midpoints.

The corresponding Riemann statement is recorded as SOH-RZ004:
if every non-trivial \(\xi\)-zero is independently established to be a global
RZ zero-mode, the triad forces \(\Re\rho=1/2\). That incoming binding remains
separate from the proved triad and therefore does not by itself alter the
external RH status.

\section*{Surviving direct G024 frontier}
The direct external criterion remains
\[
\boxed{\widehat D_y(x)>0\quad\forall x\in\mathbb R,\quad0<|y|<\frac12,}
\]
equivalently
\[
\boxed{\Re\!\left(f'(x+iy)^2-f(x+iy)f''(x+iy)\right)>0.}
\]
This is the surviving one-dimensional Jensen--Wiener frontier. It is not supplied by the disproved complete-monotonicity premise.

\section*{Zero critical relational axis and the minimum Occam channel}

Chapter~\ref{chap:zero-critical-relational-axis} compresses the relational
programme to three exact descriptions of one zero locus:
\[
J(\sigma)=\sigma
\iff
\ln2-H_2(\sigma)=0
\iff
\Delta_{\rm RC}(\Omega(s))=0
\iff
\Re s=\frac12.
\]
The first route is exchange-symmetric relational geometry, the second is
Shannon/KL information geometry, and the third is the exact projective
reciprocal--conjugation defect.  Together with the positive relational
action-defect, they identify the half-axis as the unique relational
zero-defect barycenter.

The formal Lean endgame reduces the RH channel to one incoming analytic
certificate.  If an independently constructed energy $E$ and a constant
$c>0$ satisfy, on every non-trivial zeta zero $\rho$,
\[
E(\rho)=0,
\qquad
c\,\Delta_{\rm RC}(\Omega(\rho))\le E(\rho),
\]
then
\[
\boxed{\mathrm{RH}}
\]
follows.  This is the theorem
\texttt{riemannHypothesis\_of\_coercive\_zero\_energy}.

Hence the enlarged Occam system obtained by adjoining canonical zero-energy
coercivity has a closed axiomatic RH proof channel:
\[
\boxed{\mathsf O\vdash\mathrm{RH}.}
\]
After that axiom no additional conjectural edge remains.

The same formal layer also proves that zero half-axis defect on every
non-trivial zeta zero is equivalent to RH.  Therefore the axiomatic closure
must not be confused with an independent proof in ordinary mathematics:
the canonical energy/coercivity input has to be derived from already proved,
non-RH-equivalent analytic structure before the repository can set
\texttt{proof\_of\_rh=true}.

\section*{Open problems}
\begin{enumerate}
\item \textbf{Reciprocal--conjugation zero-orbit collapse.} Prove non-circularly
\[
X(u)=0\Longrightarrow u^{-1}=\bar u,
\]
equivalently $\Delta_{\mathrm{RC}}(u)=0$.

\item \textbf{External Fourier positivity.} Prove $\widehat D_y(x)>0$ over the full frequency line and strip.

\item \textbf{SOH-G003.} Prove every zero of $F$ is real.

\item \textbf{Actual PF$_3$ and PF$_\infty$.} Close the quotient-coefficient hierarchy rather than an abstract surrogate.

\item \textbf{SOH-C005.} Prove full admissible Weil positivity without importing an RH-equivalent premise.

\item \textbf{Derivative-failure localization.} As a diagnostic question only, determine the first order $m\ge3$ and region where the G024 signed hierarchy fails. This is no longer a route to full complete monotonicity.
\end{enumerate}

\section*{Final proof boundary}
The exact projective orbit formulation, quotient real-rootedness criterion, external Fourier criterion, and classical Weil/Li criteria meet at an RH-equivalent boundary. Equivalence among them does not supply a proof. A valid closure requires a new incoming theorem from already proved structure.

Therefore the publication status is
\[
\boxed{\text{G024 full complete monotonicity: CLOSED ROUTE / NO-GO}},
\]
\[
\boxed{\text{orbit collapse, }\widehat D_y>0,\ \mathrm{SOH\!\!-G003},\ \mathrm{SOH\!\!-C005},\ \mathrm{PF}_3,\ \mathrm{PF}_\infty:\ OPEN},
\]
\[
\boxed{\mathrm{Riemann\ Hypothesis}:\ OPEN}.
\]

The Riemann Hypothesis remains OPEN in Version 0.11.
